% main.tex — Relatorio Projeto 2: Sensor Ultrasonico
% Lab de Sistemas Embarcados (EA801) — FEEC/UNICAMP
% Overleaf-ready: pdflatex ou lualatex

\documentclass[12pt, a4paper]{article}

% ── Pacotes ──────────────────────────────────────────────
\usepackage[utf8]{inputenc}
\usepackage[T1]{fontenc}
\usepackage[brazilian]{babel}
\usepackage{geometry}
\geometry{margin=2.5cm}
\usepackage{graphicx}
\usepackage{float}
\usepackage{booktabs}
\usepackage{enumitem}
\usepackage{hyperref}
\usepackage{xcolor}
\usepackage{listings}
\usepackage{caption}
\usepackage{subcaption}
\usepackage{amsmath}
\usepackage{amssymb}
\usepackage{fancyhdr}

% ── Cores e estilo ───────────────────────────────────────
\definecolor{codecomment}{rgb}{0.4,0.4,0.4}
\definecolor{codestring}{rgb}{0.2,0.5,0.2}
\definecolor{codekeyword}{rgb}{0.1,0.1,0.6}
\definecolor{codebackground}{rgb}{0.97,0.97,0.97}

\lstset{
  language=Python,
  basicstyle=\ttfamily\small,
  keywordstyle=\color{codekeyword}\bfseries,
  commentstyle=\color{codecomment}\itshape,
  stringstyle=\color{codestring},
  backgroundcolor=\color{codebackground},
  frame=single,
  breaklines=true,
  numbers=left,
  numberstyle=\tiny\color{gray},
  captionpos=b,
  tabsize=4,
}

\hypersetup{
  colorlinks=true,
  linkcolor=blue!60!black,
  urlcolor=blue!70!black,
  citecolor=green!50!black,
}

% ── Cabecalho ────────────────────────────────────────────
\pagestyle{fancy}
\fancyhf{}
\fancyhead[L]{\small EA801 --- Lab de Sistemas Embarcados}
\fancyhead[R]{\small FEEC/UNICAMP}
\fancyfoot[C]{\thepage}
\renewcommand{\headrulewidth}{0.4pt}

% ── Dados do projeto ────────────────────────────────────
\newcommand{\projetotitulo}{Projeto 2 --- Sensor Ultrassônico com Alerta WhatsApp}
\newcommand{\disciplina}{EA801 --- Laboratório de Sistemas Embarcados}
\newcommand{\instituicao}{Faculdade de Engenharia Elétrica e de Computação --- UNICAMP}

% ═════════════════════════════════════════════════════════
\begin{document}

% ── Capa ─────────────────────────────────────────────────
\begin{titlepage}
  \centering
  \vspace*{2cm}

  {\Large \textbf{\instituicao}}\\[0.5cm]
  {\large \disciplina}\\[3cm]

  {\LARGE \textbf{\projetotitulo}}\\[3cm]

  \begin{tabular}{ll}
    \toprule
    \textbf{Integrante} & \textbf{RA} \\
    \midrule
    Caio Cezar Pagliarani & 248118 \\
    Vitor Gonçalves Costa & 240043 \\
    Lucas Bernardino de Oliveira & 182199 \\
    \bottomrule
  \end{tabular}

  \vfill
  {\large Campinas, 08 de abril de 2026}
\end{titlepage}

% ── Sumario ──────────────────────────────────────────────
\tableofcontents
\newpage

% ═════════════════════════════════════════════════════════
% ETAPA 1: DEFINICAO DO PROBLEMA
% ═════════════════════════════════════════════════════════
\section{Definição do Problema}
\label{sec:problema}

\subsection{Contexto}

Este projeto é a segunda etapa de uma sequência de três projetos na disciplina EA801, com complexidade incremental:

\begin{enumerate}
  \item \textbf{Projeto 1 --- BitTeclado}: teclado musical interativo (base)
  \item \textbf{Projeto 2 --- Sensor Ultrassônico}: integração de sensor de distância com alertas remotos (este projeto)
  \item \textbf{Projeto 3 --- Bengala Assistiva}: adaptação para auxílio a pessoas com deficiência visual
\end{enumerate}

\subsection{Problema}

Monitorar a presença de objetos ou pessoas em uma zona de proximidade e notificar remotamente o usuário quando algo é detectado abaixo de uma distância crítica. O sistema deve fornecer feedback multissensorial (visual + sonoro) em tempo real e enviar alertas via WhatsApp.

\subsection{Escopo}

O sistema é um \textbf{monitor de proximidade embarcado} que:

\begin{itemize}
  \item Mede distância continuamente com sensor ultrassônico HC-SR04
  \item Fornece feedback visual graduado (OLED + NeoPixel 5$\times$5 + LED RGB)
  \item Fornece feedback sonoro proporcional à distância (buzzer, estilo sensor de ré)
  \item Envia alerta via WhatsApp quando objeto entra na zona crítica ($<$20\,cm)
  \item Conecta-se à internet via módulo ESP8266 (Wi-Fi)
\end{itemize}


% ═════════════════════════════════════════════════════════
% ETAPA 2: REQUISITOS
% ═════════════════════════════════════════════════════════
\section{Requisitos}
\label{sec:requisitos}

\subsection{Requisitos Funcionais}

\begin{table}[H]
\centering
\caption{Requisitos funcionais do sistema.}
\begin{tabular}{clp{8cm}}
  \toprule
  \textbf{RF} & \textbf{Descrição} & \textbf{Critério de Aceitação} \\
  \midrule
  RF01 & Medir distância & Leitura válida entre 2 e 400\,cm com resolução de 0,1\,cm \\
  RF02 & Feedback sonoro & Beep com frequência proporcional à distância (contínuo $<$10\,cm, silêncio $>$100\,cm) \\
  RF03 & Feedback visual (NeoPixel) & Barra de nível com cores por faixa de distância \\
  RF04 & Feedback visual (LED RGB) & Cor muda conforme faixa: verde, amarelo, vermelho \\
  RF05 & Display OLED & Exibe distância numérica, barra visual, status Wi-Fi e alerta \\
  RF06 & Conexão Wi-Fi & Conecta via ESP8266 usando comandos AT pela UART \\
  RF07 & Alerta WhatsApp & Envia mensagem via CallMeBot API quando distância $<$ 20\,cm \\
  RF08 & Cooldown de alerta & Mínimo de 30\,s entre alertas consecutivos \\
  \bottomrule
\end{tabular}
\end{table}

\subsection{Requisitos Não-Funcionais}

\begin{table}[H]
\centering
\caption{Requisitos não-funcionais do sistema.}
\begin{tabular}{clp{8cm}}
  \toprule
  \textbf{RNF} & \textbf{Descrição} & \textbf{Critério} \\
  \midrule
  RNF01 & Tempo de resposta & Leitura + feedback em $<$50\,ms por ciclo \\
  RNF02 & Plataforma & BitDogLab V7 (RP2040) com MicroPython \\
  RNF03 & Consumo & Alimentação via USB (5\,V) \\
  RNF04 & Modularidade & Código organizado em módulos independentes \\
  RNF05 & Custo & Apenas componentes já disponíveis na BitDogLab + HC-SR04 + ESP8266 \\
  \bottomrule
\end{tabular}
\end{table}


% ═════════════════════════════════════════════════════════
% ETAPA 3: CONCEITO PRELIMINAR
% ═════════════════════════════════════════════════════════
\section{Conceito Preliminar}
\label{sec:conceito}

\subsection{Abordagem de Design}

O projeto adotou a abordagem de \textbf{refinamentos sucessivos} (prototipagem incremental):

\begin{enumerate}
  \item Primeiro, validar o sensor HC-SR04 isoladamente (\texttt{teste\_sensor.py})
  \item Depois, estabelecer comunicação com o ESP8266 (\texttt{teste\_esp.py})
  \item Em seguida, integrar feedback visual e sonoro
  \item Por fim, adicionar o alerta WhatsApp via CallMeBot
\end{enumerate}

\subsection{Solução Proposta}

O RP2040 (BitDogLab) atua como controlador central, lendo o sensor a cada 50\,ms e atualizando todos os periféricos de saída. A comunicação com a internet é delegada ao ESP8266 via UART, usando comandos AT para conexão Wi-Fi e requisições HTTP.

O alerta WhatsApp utiliza a API gratuita \textbf{CallMeBot}, que permite enviar mensagens para WhatsApp via requisição HTTP GET simples.


% ═════════════════════════════════════════════════════════
% ETAPA 4: CRONOGRAMA
% ═════════════════════════════════════════════════════════
\section{Cronograma}
\label{sec:cronograma}

Este projeto é uma continuação direta do Projeto~1 (BitTeclado, entregue em 22/03). A ideia do sensor ultrassônico já havia sido concebida durante o desenvolvimento do BitTeclado, conforme descrito na seção de trabalhos futuros daquele projeto. Isso permitiu iniciar a pesquisa e aquisição de componentes (HC-SR04 e ESP8266) em paralelo à finalização do Projeto~1.

\subsection{Diagrama Gantt}

\begin{figure}[H]
\centering
\begin{verbatim}
              19/03  22/03  25/03  28/03  31/03  03/04  06/04  08/04
              QUA    SAB    TER    SEX    SEG    QUI    DOM    TER
--------------------------------------------------------------------
Fase 0: Planejamento (overlap com P1)
  Conceito + pesquisa ####|
  Aquisicao HC-SR04        ##|
  Aquisicao ESP8266        ##|

Fase 1: Validacao de Componentes
  teste_sensor.py              ####|
  teste_esp.py (baudrate)      ########|  <- dificuldade
  Montagem fisica sensor            ####|

Fase 2: Firmware Individual
  ultrasonico.py                     ####|
  buzzer.py                          ####|
  feedback.py (NeoPixel+RGB)              ####|
  display.py                              ####|

Fase 3: Wi-Fi e Integracao
  wifi.py (AT commands)                    ############|  <- maior
  Montagem ESP8266 (hardware)              ############|  <- dificuldade
  Integracao CallMeBot                               ####|
  main.py (loop completo)                             ####|

Fase 4: Testes e Entrega
  Testes integrados na placa                               ####|
  Documentacao e relatorio                                 ########|
--------------------------------------------------------------------
\end{verbatim}
\caption{Cronograma do Projeto 2 (19/03 a 08/04).}
\end{figure}

\subsection{Divisão de Tarefas por Integrante}

\subsubsection*{Caio Cezar Pagliarani (248118) --- Arquitetura, Wi-Fi e Integração}

\begin{table}[H]
\centering
\begin{tabular}{lllc}
  \toprule
  \textbf{Tarefa} & \textbf{Início} & \textbf{Fim} & \textbf{Status} \\
  \midrule
  Conceito do projeto e pesquisa de viabilidade & 19/03 & 22/03 & OK \\
  Repositório GitHub e estrutura inicial & 22/03 & 22/03 & OK \\
  Módulo \texttt{wifi.py} (AT com ESP8266) & 28/03 & 03/04 & OK \\
  Integração \texttt{main.py} (loop principal) & 03/04 & 05/04 & OK \\
  Testes integrados e debug de timing UART & 05/04 & 06/04 & OK \\
  Relatório LaTeX e documentação final & 06/04 & 08/04 & OK \\
  \bottomrule
\end{tabular}
\end{table}

\subsubsection*{Vitor Gonçalves Costa (240043) --- Hardware e Sensor}

\begin{table}[H]
\centering
\begin{tabular}{lllc}
  \toprule
  \textbf{Tarefa} & \textbf{Início} & \textbf{Fim} & \textbf{Status} \\
  \midrule
  Aquisição e teste inicial do HC-SR04 & 22/03 & 25/03 & OK \\
  Módulo \texttt{ultrasonico.py} (driver do sensor) & 25/03 & 26/03 & OK \\
  Módulo \texttt{buzzer.py} (beep proporcional) & 26/03 & 28/03 & OK \\
  Montagem física do ESP8266 & 28/03 & 03/04 & OK \\
  Diagnóstico de mau contato e remontagem & 31/03 & 03/04 & OK \\
  Documentação de hardware e pinout & 06/04 & 07/04 & OK \\
  \bottomrule
\end{tabular}
\end{table}

\subsubsection*{Lucas Bernardino de Oliveira (182199) --- Display e Feedback Visual}

\begin{table}[H]
\centering
\begin{tabular}{lllc}
  \toprule
  \textbf{Tarefa} & \textbf{Início} & \textbf{Fim} & \textbf{Status} \\
  \midrule
  Pesquisa NeoPixel e mapeamento da matriz 5$\times$5 & 22/03 & 25/03 & OK \\
  Módulo \texttt{feedback.py} (NeoPixel + LED RGB) & 25/03 & 28/03 & OK \\
  Módulo \texttt{display.py} (painel OLED) & 28/03 & 31/03 & OK \\
  Script \texttt{teste\_esp.py} (varredura de baudrate) & 28/03 & 29/03 & OK \\
  Testes unitários dos módulos visuais & 31/03 & 01/04 & OK \\
  Apresentação e revisão final & 06/04 & 08/04 & OK \\
  \bottomrule
\end{tabular}
\end{table}

\subsection{Marcos (Milestones)}

\begin{table}[H]
\centering
\begin{tabular}{lll}
  \toprule
  \textbf{Data} & \textbf{Marco} & \textbf{Critério de Aceitação} \\
  \midrule
  22/03 & M0 --- Entrega P1 + início P2 & P1 entregue, repo P2 criado \\
  25/03 & M1 --- Sensor validado & HC-SR04 medindo corretamente \\
  28/03 & M2 --- Módulos individuais prontos & buzzer, display, feedback OK \\
  03/04 & M3 --- Wi-Fi funcional & ESP8266 conectando e enviando HTTP \\
  05/04 & M4 --- Integração completa & main.py com todos os módulos \\
  08/04 & M5 --- Entrega final & Relatório + código + apresentação \\
  \bottomrule
\end{tabular}
\end{table}

\subsection{Dificuldades no Cronograma}

A Fase~3 (Wi-Fi e Integração) consumiu significativamente mais tempo que o planejado, por três razões principais:

\begin{enumerate}
  \item \textbf{Mau contato nos jumpers do ESP8266}: a conexão física entre a BitDogLab e o módulo ESP-01 apresentou intermitência. Jumpers soltos e protoboard com contatos desgastados causaram falhas de comunicação UART que inicialmente foram confundidas com problemas de software.
  \item \textbf{Baudrate e timing dos comandos AT}: foi necessário criar o script \texttt{teste\_esp.py} para varrer todos os baudrates possíveis. Mesmo após encontrar o baudrate correto (115200), os tempos de espera dos comandos AT precisaram de ajuste empírico.
  \item \textbf{Erros humanos na montagem}: inversão de pinos TX/RX, alimentação incorreta (5\,V em vez de 3.3\,V no ESP8266), e fios conectados nos GPIOs errados atrasaram a validação por vários dias.
\end{enumerate}


% ═════════════════════════════════════════════════════════
% ETAPA 5: ARQUITETURA
% ═════════════════════════════════════════════════════════
\section{Arquitetura do Sistema}
\label{sec:arquitetura}

\subsection{Diagrama de Blocos}

\begin{figure}[H]
\centering
\begin{verbatim}
  +-------------+                    +--------------+
  |  HC-SR04    |--GPIO8/9---------->|              |--GPIO21--> Buzzer PWM
  |  (sensor)   |                    |              |
  +-------------+                    |   RP2040     |
                                     |  (BitDogLab) |--GPIO7---> NeoPixel 5x5
  +-------------+                    |              |
  |  ESP8266    |<-UART(GP0/1)------>|              |--GP13/11/12-> LED RGB
  |  (Wi-Fi)    |                    |              |
  +-------------+                    |              |--I2C(GP2/3)-> OLED SSD1306
       |                             +--------------+
       | HTTP
       v
  +-------------+
  | CallMeBot   |
  | (WhatsApp)  |
  +-------------+
\end{verbatim}
\caption{Diagrama de blocos do sistema.}
\label{fig:diagrama-texto}
\end{figure}

\subsection{Módulos de Software}

O firmware é organizado em 6 módulos independentes:

\begin{table}[H]
\centering
\caption{Módulos do firmware e suas responsabilidades.}
\begin{tabular}{lp{9cm}}
  \toprule
  \textbf{Módulo} & \textbf{Responsabilidade} \\
  \midrule
  \texttt{main.py} & Loop principal: orquestra leitura, feedback e alertas \\
  \texttt{ultrasonico.py} & Driver do HC-SR04: pulso trigger e medição do echo \\
  \texttt{buzzer.py} & Controle PWM do buzzer com beep proporcional à distância \\
  \texttt{display.py} & Painel OLED: distância, barra visual, status Wi-Fi/alerta \\
  \texttt{feedback.py} & Controle dos NeoPixels 5$\times$5 e LED RGB por faixas de cor \\
  \texttt{wifi.py} & Comunicação UART com ESP8266 (AT commands) + CallMeBot \\
  \bottomrule
\end{tabular}
\end{table}

\subsection{Fluxo de Execução}

O loop principal executa a cada 50\,ms (\texttt{INTERVALO\_LEITURA\_MS}) na seguinte ordem:

\begin{enumerate}
  \item \textbf{Medir} distância via HC-SR04 (pulso de 10\,$\mu$s no trigger)
  \item \textbf{Buzzer}: ajustar beep conforme distância
  \item \textbf{NeoPixel + LED RGB}: atualizar cores e barra de nível
  \item \textbf{Verificar alerta}: se distância $<$ 20\,cm e cooldown expirou, enviar WhatsApp
  \item \textbf{OLED}: atualizar display com dados atuais
\end{enumerate}


% ═════════════════════════════════════════════════════════
% ETAPA 6: ESPECIFICACOES TECNICAS
% ═════════════════════════════════════════════════════════
\section{Especificações Técnicas}
\label{sec:especificacoes}

\subsection{Lista de Materiais (BOM)}

\begin{table}[H]
\centering
\caption{Bill of Materials.}
\begin{tabular}{llcl}
  \toprule
  \textbf{Componente} & \textbf{Modelo} & \textbf{Qtd} & \textbf{Observação} \\
  \midrule
  Placa de desenvolvimento & BitDogLab V7 (RP2040) & 1 & Inclui OLED, NeoPixel, RGB, buzzer \\
  Sensor ultrassônico & HC-SR04 & 1 & Alcance 2--400\,cm \\
  Módulo Wi-Fi & ESP8266 (ESP-01) & 1 & Firmware AT, 3.3\,V \\
  Jumpers / fios & --- & $\sim$6 & Conexão sensor e ESP \\
  \bottomrule
\end{tabular}
\end{table}

\subsection{Mapeamento de Pinos}

\begin{table}[H]
\centering
\caption{Mapeamento de GPIOs utilizados.}
\begin{tabular}{llll}
  \toprule
  \textbf{Periférico} & \textbf{Pino} & \textbf{Direção} & \textbf{Interface} \\
  \midrule
  HC-SR04 Trigger & GPIO8 & Saída & Digital \\
  HC-SR04 Echo    & GPIO9 & Entrada & Digital \\
  Buzzer          & GPIO21 & Saída & PWM \\
  NeoPixel 5$\times$5 & GPIO7 & Saída & WS2812 \\
  LED Vermelho    & GPIO13 & Saída & PWM \\
  LED Verde       & GPIO11 & Saída & PWM \\
  LED Azul        & GPIO12 & Saída & PWM \\
  OLED SDA        & GPIO2 & Bidirecional & I$^2$C \\
  OLED SCL        & GPIO3 & Saída & I$^2$C \\
  ESP8266 TX      & GPIO0 & Saída & UART0 \\
  ESP8266 RX      & GPIO1 & Entrada & UART0 \\
  \bottomrule
\end{tabular}
\end{table}

\subsection{Princípio de Funcionamento do HC-SR04}

O sensor ultrassônico opera por tempo de voo (\textit{time-of-flight}):

\begin{enumerate}
  \item O microcontrolador envia um pulso HIGH de 10\,$\mu$s no pino Trigger
  \item O sensor emite 8 pulsos ultrassônicos a 40\,kHz
  \item O pino Echo fica HIGH pelo tempo que o som leva para ir e voltar
  \item A distância é calculada por:
\end{enumerate}

\begin{equation}
  d = \frac{t_{\text{echo}}}{2} \times v_{\text{som}} = \frac{t_{\text{echo}} \; [\mu s]}{2 \times 29{,}15} \; [\text{cm}]
\end{equation}

onde $v_{\text{som}} \approx 343$\,m/s $= 29{,}15$\,$\mu$s/cm.


% ═════════════════════════════════════════════════════════
% ETAPA 7: PROTOTIPAGEM E TESTES
% ═════════════════════════════════════════════════════════
\section{Prototipagem e Testes}
\label{sec:prototipagem}

\subsection{Teste 1: Sensor HC-SR04}

O primeiro teste validou a leitura do sensor isoladamente, usando o script \texttt{teste\_sensor.py}:

\begin{lstlisting}[caption={Trecho do teste standalone do HC-SR04.}]
trig = Pin(8, Pin.OUT)
echo = Pin(9, Pin.IN)

# Pulso de trigger: 10us HIGH
trig.value(1)
time.sleep_us(10)
trig.value(0)

# Medir duracao do echo (timeout 30ms)
duracao = time_pulse_us(echo, 1, 30000)
dist = (duracao / 2.0) / 29.15
print("Distancia: {:.1f} cm".format(dist))
\end{lstlisting}

\subsection{Teste 2: Comunicação ESP8266}

O script \texttt{teste\_esp.py} testou a comunicação UART com o ESP8266 varrendo múltiplos baudrates (9600, 19200, 38400, 57600, 115200, 74880) até encontrar resposta ``OK'' ao comando AT.

O baudrate correto encontrado foi \textbf{115200\,bps}.

\subsection{Teste 3: Integração Completa}

\begin{itemize}
  \item \textbf{Feedback sonoro}: validado nas 4 faixas de distância (contínuo, 100\,ms, 300\,ms, 600\,ms)
  \item \textbf{NeoPixel}: barra de nível funcionando com cores corretas por faixa
  \item \textbf{LED RGB}: transição de cores conforme proximidade
  \item \textbf{OLED}: atualização sem flickering (cache de estado anterior)
  \item \textbf{WhatsApp}: mensagem recebida em $\sim$3--5\,s após detecção
\end{itemize}


% ═════════════════════════════════════════════════════════
% ETAPA 8: VALIDACAO
% ═════════════════════════════════════════════════════════
\section{Validação}
\label{sec:validacao}

\subsection{Checklist de Requisitos}

\begin{table}[H]
\centering
\caption{Validação dos requisitos funcionais.}
\begin{tabular}{clc}
  \toprule
  \textbf{Req.} & \textbf{Descrição} & \textbf{Atendido?} \\
  \midrule
  RF01 & Medição de distância 2--400\,cm & Sim \\
  RF02 & Beep proporcional à distância & Sim \\
  RF03 & Barra NeoPixel por faixa de cor & Sim \\
  RF04 & LED RGB por faixa de distância & Sim \\
  RF05 & Display OLED com info completa & Sim \\
  RF06 & Conexão Wi-Fi via ESP8266 & Sim \\
  RF07 & Alerta WhatsApp via CallMeBot & Sim \\
  RF08 & Cooldown de 30\,s entre alertas & Sim \\
  \bottomrule
\end{tabular}
\end{table}


% ═════════════════════════════════════════════════════════
% CONCLUSAO
% ═════════════════════════════════════════════════════════
\section{Conclusão}
\label{sec:conclusao}

O Projeto~2 representou um salto de complexidade em relação ao BitTeclado (Projeto~1), passando de um sistema puramente local para uma aplicação conectada à internet com comunicação entre dois microcontroladores (RP2040 e ESP8266). O resultado final é um monitor de proximidade embarcado funcional, capaz de medir distância em tempo real, fornecer feedback multissensorial e enviar alertas remotos via WhatsApp.

\subsection{Principais Resultados}

\begin{itemize}
  \item Sistema de medição de distância funcional com 4 faixas de feedback sonoro (estilo sensor de ré) e visual (barra de nível NeoPixel + LED RGB com código de cores)
  \item Comunicação Wi-Fi operacional via ESP8266 com comandos AT, incluindo requisições HTTP para a API CallMeBot
  \item Alerta WhatsApp recebido em 3--5\,s após detecção de objeto na zona crítica ($<$20\,cm), com cooldown de 30\,s para evitar spam
  \item Arquitetura modular com 6 módulos independentes, reaproveitando padrões do Projeto~1 (buzzer, display OLED)
  \item Display OLED com atualização otimizada por cache de estado, eliminando flickering
\end{itemize}

\subsection{Dificuldades Encontradas}

A principal fonte de dificuldades foi a integração do módulo ESP8266, em três dimensões:

\begin{enumerate}
  \item \textbf{Hardware --- mau contato}: a conexão entre BitDogLab e ESP-01 via protoboard e jumpers apresentou intermitência frequente. Contatos desgastados na protoboard e jumpers soltos causaram falhas de comunicação UART que, inicialmente, foram diagnosticadas erroneamente como bugs de software. A solução envolveu refazer a montagem múltiplas vezes, trocar jumpers e verificar continuidade com multímetro.

  \item \textbf{Hardware --- erros humanos}: durante a montagem, ocorreram inversões de pinos TX/RX e, em uma ocasião, alimentação de 5\,V aplicada diretamente ao ESP8266 (que opera a 3.3\,V). Esses erros atrasaram a validação e exigiram cuidado redobrado na verificação do circuito antes de cada teste.

  \item \textbf{Software --- comandos AT}: a comunicação UART com o ESP8266 exigiu ajuste empírico de tempos de espera (\textit{timeouts}) para cada comando AT. O comando \texttt{AT+CWJAP} (conexão Wi-Fi) precisou de até 10\,s de espera, enquanto o envio HTTP (\texttt{AT+CIPSEND}) demandou pausas intermediárias de 500\,ms para sincronização. Foi criado o script \texttt{teste\_esp.py} com varredura automática de baudrates para diagnosticar problemas de comunicação.

  \item \textbf{Dashboard descontinuado}: uma tentativa de criar um dashboard em tempo real (Flask + Chart.js no PC) foi abandonada por instabilidade na comunicação HTTP entre o ESP8266 e o servidor local, além de complexidade excessiva para o escopo do projeto.
\end{enumerate}

\subsection{Lições Aprendidas}

\begin{itemize}
  \item \textbf{Testar hardware antes de software}: quando a comunicação falha, verificar sempre a camada física primeiro (jumpers, alimentação, pinout) antes de depurar código.
  \item \textbf{Scripts de diagnóstico são essenciais}: o \texttt{teste\_sensor.py} e o \texttt{teste\_esp.py} economizaram tempo ao isolar problemas de componentes individuais.
  \item \textbf{Modularidade facilita debug}: como cada módulo (sensor, buzzer, display, Wi-Fi) era independente, foi possível testar e validar cada um isoladamente antes da integração.
  \item \textbf{Reusar código entre projetos}: padrões do Projeto~1 (estrutura do buzzer, display OLED) foram adaptados diretamente, acelerando o desenvolvimento.
\end{itemize}

\subsection{Trabalhos Futuros}

\begin{itemize}
  \item \textbf{Projeto 3 --- Bengala assistiva}: adaptação do sistema de medição ultrassônica e feedback sonoro para uma bengala eletrônica destinada a pessoas com deficiência visual, conforme planejado desde a concepção do Projeto~1.
  \item Substituição do ESP-01 por ESP32 integrado para eliminar problemas de comunicação UART e simplificar a montagem.
  \item Implementação de filtro de média móvel nas leituras do HC-SR04 para reduzir ruído em ambientes com superfícies irregulares.
\end{itemize}


% ═════════════════════════════════════════════════════════
% REFERENCIAS
% ═════════════════════════════════════════════════════════
\section*{Referências}
\addcontentsline{toc}{section}{Referências}

\begin{enumerate}[label={[\arabic*]}]
  \item ARES, V. E. \textit{Metodologia Projeto v5}. FEEC/UNICAMP, 2024.
  \item CUGNASCA, C. E. \textit{Projeto de Sistemas Embarcados}. Escola Politécnica, USP, 2018.
  \item Elecfreaks. \textit{HC-SR04 Ultrasonic Sensor Datasheet}. Disponível em: \url{https://cdn.sparkfun.com/datasheets/Sensors/Proximity/HCSR04.pdf}
  \item Espressif Systems. \textit{ESP8266 AT Instruction Set}. Disponível em: \url{https://www.espressif.com/sites/default/files/documentation/4a-esp8266_at_instruction_set_en.pdf}
  \item CallMeBot. \textit{WhatsApp API Documentation}. Disponível em: \url{https://www.callmebot.com/blog/free-api-whatsapp-messages/}
\end{enumerate}

\end{document}
